% Chapter 8: Computational Complexity
% Part III: Difficulty

\chapter{Computational Complexity}
\label{ch:complexity}

%======================================================================
% SECTION 1: THE QUESTION
%======================================================================
\section{The Question: What Problem Was Humanity Trying to Solve?}
\label{sec:complexity:question}

Turing (1936) answered ``What is computable?'' But not all computable problems are equal. Some take seconds; others take longer than the age of the universe. The question that birthed complexity theory:

What makes a problem \emph{inherently} difficult? Can we classify problems by their intrinsic computational cost, independent of algorithms or hardware?

Is there a fundamental difference between problems solvable in polynomial time and those requiring exponential time? Can we prove that certain problems \emph{require} exponential resources — or is every problem in EXP actually in P, waiting for a clever algorithm?

Before complexity theory, ``difficulty'' was not a mathematical concept. A sorting algorithm might be ``slow'' for large inputs, but the problem of sorting itself was not considered intrinsically hard or easy. The radical idea: difficulty is a property of the \emph{problem}, not the \emph{algorithm} — and it can be measured, classified, and proven.

The distinction between polynomial and exponential growth is stark. A problem solvable in $n^3$ steps on a $10^9$ ops/sec machine can handle inputs of size $10^6$ within seconds. An exponential $2^n$ algorithm struggles at $n = 50$. This gap defines the boundary between the feasible and the intractable.

Even after defining polynomial time as ``efficient,'' we face a deeper question: can we determine, for any given problem, which class it belongs to? The P vs NP question — are efficient verification and efficient solution the same? — remains open after 50 years and \$1M Clay Prize. The entire field orbits this single question.

Kolmogorov complexity measures the \emph{descriptive} difficulty of individual strings. Computational complexity measures the \emph{computational} difficulty of \emph{sets} of strings (languages). Both confront the same question: what does it mean for something to be ``hard?'' Chapter 7 defines randomness via incompressibility; this chapter defines intractability via resource lower bounds.

%======================================================================
% SECTION 2: WHY IT SEEMED IMPOSSIBLE
%======================================================================
\section{Why It Seemed Impossible: What Barriers Blocked Progress}
\label{sec:complexity:impossible}

Before the 1960s, ``difficulty'' was not a mathematical concept:

\begin{enumerate}
\item \textbf{No Formal Model of Cost:} Turing machines compute functions; they don't measure \emph{how long} they take. Time/space as complexity measures required a new layer of theory.
\item \textbf{Algorithm-Centric View:} Difficulty was attributed to \emph{algorithms} (``this sorting algorithm is slow''), not \emph{problems} (``sorting is inherently hard''). The idea that a problem has an \emph{intrinsic} difficulty, independent of the algorithm used, was radical.
\item \textbf{No Lower Bounds:} Proving that \emph{no} algorithm can solve a problem faster than $f(n)$ seemed impossible. You'd have to reason about \emph{all possible algorithms} — including ones not yet invented.
\item \textbf{The ``Clever Programmer'' Hope:} For every ``hard'' problem, experts believed a clever algorithm would eventually make it easy. The concept of \emph{inherent intractability} was foreign.
\end{enumerate}

Before the 1960s, the only lower bounds were trivial: any algorithm solving a problem that depends on all $n$ input bits must read them all ($\Omega(n)$). Subroutines like comparison-based sorting had $\Omega(n \log n)$ lower bounds but these were model-specific (comparison trees), not about TMs. The idea of a \emph{universal} lower bound — one that holds for every possible algorithm in any reasonable model — was considered hopeless.

Proving upper bounds is easy: exhibit an algorithm. Proving lower bounds requires showing \emph{no algorithm exists} — a universal quantification over all possible algorithms. This is why we have P $\neq$ NP (conjectured) but almost no unconditional lower bounds beyond trivial ones (reading input takes $\Omega(n)$). The only unconditional separations we have are via diagonalization (hierarchy theorems), which work by constructing problems with built-in self-reference — not by analyzing natural problems.

Baker, Gill, and Solovay (1975) showed that there exist oracles $A$ and $B$ such that $\mathsf{P}^A = \mathsf{NP}^A$ and $\mathsf{P}^B \neq \mathsf{NP}^B$. This means any proof resolving P vs NP must be \emph{non-relativizing} — it cannot treat all TMs uniformly as black boxes. This eliminated diagonalization (which relativizes) as a viable approach. Later, Aaronson and Wigderson (2008) identified an even broader class of techniques (algebrization) that also cannot resolve the question.

%======================================================================
% SECTION 3: HISTORICAL CONTEXT
%======================================================================
\section{Historical Context: What Was Known at the Time}
\label{sec:complexity:context}

\begin{itemize}
\item \textbf{1960:} Rabin: ``Degree of difficulty of computing a function.''
\item \textbf{1965:} Hartmanis \& Stearns: ``On the Computational Complexity of Algorithms'' — defined time/space complexity classes, hierarchy theorems, founded the field.
\item \textbf{1965:} \cite{edmonds1965}, Edmonds: ``Paths, Trees, and Flowers'' — proposed \emph{polynomial time} as the formal definition of ``tractable'' (``good algorithms'').
\item \textbf{1967:} \cite{cobham1965}, Cobham: ``The Intrinsic Computational Difficulty of Functions'' — independently proposed polynomial time = feasible.
\item \textbf{1971:} Cook: ``The Complexity of Theorem-Proving Procedures'' — NP-completeness, SAT is NP-complete.
\item \textbf{1972:} Karp: ``Reducibility Among Combinatorial Problems'' — 21 NP-complete problems.
\item \textbf{1973:} Levin (USSR): Universal search problems — independent discovery of NP-completeness.
\end{itemize}

Born in Latvia, Hartmanis experienced World War II as a refugee, fleeing the Soviet occupation. He earned his PhD from Caltech in 1955 and joined Cornell, where he collaborated with Richard Stearns. Their 1965 paper founded computational complexity theory. Hartmanis later served as Cornell's department chair and received the Turing Award in 1993 (with Stearns). His work bridged the gap between recursive function theory and practical computation.

Stearns earned his PhD from Princeton (1961) working on game theory and sequential machines. At GE Research Lab, he met Hartmanis at a conference and they began a collaboration that produced the hierarchy theorems. Stearns also contributed to automata theory, distributed computing (the Byzantine Generals problem), and received the Turing Award (1993). His career exemplifies how theoretical insight can emerge from industrial research laboratories.

Edmonds revolutionized combinatorial optimization. His 1965 paper ``Paths, Trees, and Flowers'' gave the first polynomial-time algorithm for maximum matching in general graphs (the ``blossom algorithm''). In it, he explicitly defined polynomial time as the criterion for an algorithm to be ``good.'' This was a deliberate choice: ``For practical purposes, the difference between algebraic and exponential order is more crucial than the difference between finite and non-finite.'' He also proved that matroid intersection is in P and conjectured that P $\neq$ NP.

Cobham independently formulated the polynomial-time thesis in his 1965 paper ``The Intrinsic Computational Difficulty of Functions.'' Writing from IBM's Watson Research Center, he proposed that problems solvable in polynomial time form the natural boundary of tractability — now called ``Cobham's thesis.'' His work was less algorithmically oriented than Edmonds' but equally influential in formalizing the notion of efficient computation.

Key precursors:
\begin{itemize}
\item \textbf{1936:} Turing machines (the model).
\item \textbf{1956:} G\"odel's letter to von Neumann — first mention of ``polynomial time'' as a barrier. G\"odel asked whether a problem (SAT) could be solved in $kn$ or $kn^2$ steps, presaging the P vs NP question by 15 years.
\item \textbf{1960s:} Early attempts at lower bounds (switching circuits, formula size).
\item \textbf{1959--1962:} Rabin, Yamada, and others studied time-bounded computation but lacked the key insight of polynomial time as a robust class.
\end{itemize}

Hartmanis-Stearns used \emph{multitape Turing machines} as the cost model. The \emph{invariance thesis} (extended Church-Turing): any ``reasonable'' model of computation can be simulated on a TM with polynomial overhead. This makes P, NP, EXP \emph{model-independent} classes. Chapter 1 asked ``what is computation?'' — this chapter asks ``what is the \emph{cost} of computation?''

%======================================================================
% SECTION 4: THE MENTAL LEAP
%======================================================================
\section{The Mental Leap: The Creative Insight That Changed Everything}
\label{sec:complexity:leap}

\emph{Polynomial time} ($n^{O(1)}$) is the correct mathematical formalization of ``tractable,'' ``feasible,'' ``efficiently computable.''

This is not obvious: why not $n \log n$? Why not $2^{\sqrt{n}}$? Why not $n^{100}$? The answer: \textbf{polynomials are closed under composition}. If subroutine A runs in $p(n)$ and calls subroutine B running in $q(n)$, the total is $p(q(n))$ — still polynomial. This \emph{modularity} makes polynomial time the only complexity class that supports abstraction and composition — the essence of engineering.

The second leap (Cook 1971, Levin 1973, Karp 1972):

\emph{Reduction} is the right tool for comparing difficulty. If problem A reduces to problem B (A $\le_p$ B), then B is ``at least as hard as'' A. This creates a \emph{partial order} on problems. The \emph{maximal} elements of NP under this order are the \emph{NP-complete} problems — the ``hardest'' problems in NP. If \emph{any} NP-complete problem is in P, then P = NP.

Together, these insights created a \emph{taxonomy of difficulty}:
\begin{itemize}
\item P: Efficiently solvable.
\item NP: Efficiently verifiable.
\item NP-complete: Hardest in NP (mutually reducible).
\item EXP: Exponential time.
\item PSPACE, L, NL, BPP, PH, \dots
\end{itemize}

\begin{example}[Why Polynomial Time Is the Right Notion of Feasibility]
Consider three algorithms for a problem:
\begin{align*}
\text{Alg 1:} &\quad n^3 \text{ steps} \\
\text{Alg 2:} &\quad 2^{\sqrt{n}} \text{ steps} \\
\text{Alg 3:} &\quad 2^n \text{ steps}
\end{align*}
At $n = 10^6$, Alg 1 takes $10^{18}$ steps ($\approx 30$ years at 1 GHz). Alg 2 takes $2^{1000} \approx 10^{301}$ steps — physically impossible. But Alg 1 can be combined with other poly-time subroutines and remain polynomial; Alg 2 cannot. The boundary is not about absolute speed but about \emph{composability}: polynomial time is the only class closed under polynomial reductions.
\end{example}

\begin{example}[A Concrete Reduction: Independent Set $\le_p$ Vertex Cover]
Given a graph $G = (V,E)$, Independent Set asks: is there a set $S \subseteq V$ of size $k$ with no edges inside? Vertex Cover asks: is there a set $C \subseteq V$ of size $k$ touching all edges? Observation: $S$ is an independent set iff $V \setminus S$ is a vertex cover. So $G$ has an independent set of size $k$ iff $G$ has a vertex cover of size $n-k$. This polynomial-time reduction shows the two problems are essentially the same — solving either solves the other.
\end{example}

%======================================================================
% SECTION 5: THE MATHEMATICS
%======================================================================
\section{The Mathematics: Rigorous Development of the Theory}
\label{sec:complexity:math}

\subsection{Complexity Classes}
\label{sec:complexity:classes}

\begin{definition}[DTIME, NTIME, DSPACE, NSPACE]
$\mathsf{DTIME}(t(n))$ = languages decided by deterministic TM in $O(t(n))$ time.
$\mathsf{NTIME}(t(n))$ = languages decided by non-deterministic TM in $O(t(n))$ time.
$\mathsf{DSPACE}(s(n))$, $\mathsf{NSPACE}(s(n))$ analogously for space.
\end{definition}

\begin{definition}[Standard Classes]
\begin{align*}
\mathsf{P} &= \bigcup_{k} \mathsf{DTIME}(n^k) \\
\mathsf{NP} &= \bigcup_{k} \mathsf{NTIME}(n^k) \\
\mathsf{EXP} &= \bigcup_{k} \mathsf{DTIME}(2^{n^k}) \\
\mathsf{NEXP} &= \bigcup_{k} \mathsf{NTIME}(2^{n^k}) \\
\mathsf{PSPACE} &= \bigcup_{k} \mathsf{DSPACE}(n^k) \\
\mathsf{L} &= \mathsf{DSPACE}(\log n) \\
\mathsf{NL} &= \mathsf{NSPACE}(\log n) \\
\mathsf{BPP} &= \text{probabilistic poly-time, error } \le 1/3 \\
\mathsf{coNP} &= \{ L \mid \overline{L} \in \mathsf{NP} \} \\
\mathsf{PH} &= \bigcup_{i \ge 0} \Sigma_i^{\mathsf{P}} \quad \text{(Polynomial Hierarchy)}
\end{align*}
\end{definition}

\begin{definition}[coNP and Complements]
A language $L$ is in coNP if its complement $\overline{L}$ is in NP. Equivalently: ``no'' instances have polynomial-size certificates. The question NP = coNP is almost as deep as P = NP. If NP $\neq$ coNP, then P $\neq$ NP. (Because P is closed under complement, so P = NP would imply NP = coNP.)
\end{definition}

\begin{definition}[Polynomial Hierarchy (Stockmeyer 1976)]
Define $\Sigma_0^{\mathsf{P}} = \Pi_0^{\mathsf{P}} = \mathsf{P}$. For $i \ge 0$:
\begin{align*}
\Sigma_{i+1}^{\mathsf{P}} &= \mathsf{NP}^{\Sigma_i^{\mathsf{P}}} \\
\Pi_{i+1}^{\mathsf{P}} &= \mathsf{coNP}^{\Sigma_i^{\mathsf{P}}}
\end{align*}
The Polynomial Hierarchy $\mathsf{PH} = \bigcup_{i \ge 0} \Sigma_i^{\mathsf{P}}$. If any level collapses ($\Sigma_i^{\mathsf{P}} = \Pi_i^{\mathsf{P}}$ or $\Sigma_i^{\mathsf{P}} = \Sigma_{i+1}^{\mathsf{P}}$), the whole hierarchy collapses to that level. P = NP would collapse PH to P.
\end{definition}

\subsection{Hierarchy Theorems}
\label{sec:complexity:hierarchy}

The hierarchy theorems were the first — and remain among the only — unconditional separations of complexity classes. They use \emph{diagonalization}, a technique inherited from computability theory (Chapter 2).

\begin{theorem}[Time Hierarchy Theorem (Hartmanis-Stearns 1965)]
If $t_1(n) \log t_1(n) = o(t_2(n))$, then $\mathsf{DTIME}(t_1(n)) \subsetneq \mathsf{DTIME}(t_2(n))$.
\end{theorem}

\begin{proof}
We construct a Turing machine $D$ that decides a language $L(D)$ in $\mathsf{DTIME}(t_2)$ but not in $\mathsf{DTIME}(t_1)$. 

$D$ operates as follows on input $x$ of length $n = |x|$:
\begin{enumerate}
\item Interpret $x$ as $\langle M \rangle 10^k$ where $M$ is a TM description (if $x$ is not of this form, reject). Let $m$ be the length of $\langle M \rangle$.
\item Compute $t_2(n)$. Initialize a counter to $t_2(n)$.
\item Simulate $M$ on input $\langle M \rangle$ while decrementing the counter. If the counter reaches 0 before $M$ halts, reject.
\item If $M$ halts and accepts, \emph{reject}; if $M$ halts and rejects, \emph{accept}.
\end{enumerate}

\textbf{Correctness:} $D$ always halts within $O(t_2(n))$ steps because the counter enforces the bound. If $M$ runs in $t_1(n)$ time, then for sufficiently large $n$ (where $t_2(n) > t_1(m) \log t_1(m)$), the simulation completes, and $D$ flips $M$'s answer. Thus $L(D) \neq L(M)$ for any $M$ running in $t_1$ time. Therefore $L(D) \notin \mathsf{DTIME}(t_1)$.

\textbf{Time analysis:} Simulating $M$ for $t$ steps on a multitape TM takes $O(t \log t)$ steps (due to the overhead of simulating $M$'s tape layout and state transitions). Since $t_1(n) \log t_1(n) = o(t_2(n))$, for large $n$ we have $t_1(m) \log t_1(m) \le t_2(n)$, and $D$ runs in $O(t_2(n))$ time overall. So $L(D) \in \mathsf{DTIME}(t_2)$.
\end{proof}

\begin{corollary}
$\mathsf{P} \subsetneq \mathsf{EXP}$. (Take $t_1 = n^k$, $t_2 = 2^n$.)
\end{corollary}

\begin{corollary}
$\mathsf{DTIME}(n) \subsetneq \mathsf{DTIME}(n^2) \subsetneq \mathsf{DTIME}(n^3) \subsetneq \dots$. Time classes form a strict hierarchy, at least for functions separated by a $\log$ factor.
\end{corollary}

\begin{theorem}[Space Hierarchy Theorem]
If $s_1(n) = o(s_2(n))$, then $\mathsf{DSPACE}(s_1(n)) \subsetneq \mathsf{DSPACE}(s_2(n))$.
\end{theorem}

\begin{proof}
The construction is similar to the time hierarchy, but space is easier: we do not need the $\log$ factor. Construct a machine $D$ that on input $x = \langle M \rangle 10^k$:
\begin{enumerate}
\item Simulates $M$ on $\langle M \rangle$, keeping a counter of space used.
\item If the simulation would exceed $s_2(n)$ space, reject.
\item If $M$ halts, flip the answer.
\end{enumerate}
The simulation overhead for space is a constant factor ($M$'s tape alphabet may need encoding). Since $s_1 = o(s_2)$, for large enough $n$ the simulation fits within $s_2(n)$. Thus $L(D) \in \mathsf{DSPACE}(s_2)$ but $L(D) \notin \mathsf{DSPACE}(s_1)$.
\end{proof}

\begin{corollary}
$\mathsf{L} \subsetneq \mathsf{PSPACE}$ and $\mathsf{PSPACE} \subsetneq \mathsf{EXP}$. Hierarchy theorems give us these separations unconditionally — a rare gift in complexity theory.
\end{corollary}

\begin{theorem}[Savitch's Theorem (1970)]
$\mathsf{NSPACE}(s(n)) \subseteq \mathsf{DSPACE}(s(n)^2)$ for $s(n) \ge \log n$.
\end{theorem}

\begin{proof}
Let $N$ be an NTM using $s(n)$ space. For a fixed input $x$ of length $n$, consider $N$'s \emph{configuration graph} $G_N$: vertices are all possible configurations reachable on space $s = s(n)$, with edges representing non-deterministic transitions. Each configuration includes the tape contents, head position, and state. The number of configurations is $|Q| \cdot (n + s) \cdot |\Gamma|^{O(s)} = 2^{O(s)}$, where $Q$ are states and $\Gamma$ is the tape alphabet.

Define $\text{CANYIELD}(c_1, c_2, t)$: is there a path of length at most $2^t$ from $c_1$ to $c_2$? We compute this recursively:
\begin{align*}
\text{CANYIELD}(c_1, c_2, 0) &\iff c_1 = c_2 \text{ or } c_1 \to c_2 \text{ in one step} \\
\text{CANYIELD}(c_1, c_2, t) &\iff \exists c_m : \text{CANYIELD}(c_1, c_m, t-1) \land \text{CANYIELD}(c_m, c_2, t-1)
\end{align*}
The recursive algorithm explores all possible midpoints $c_m$. The depth of recursion is $O(\log(2^{O(s)})) = O(s)$. At each level, we store $O(s)$ bits for the current configuration. Total space: $O(s^2)$.

Finally, to decide if $N$ accepts $x$, we check $\text{CANYIELD}(c_{\text{start}}, c_{\text{accept}}, O(s))$ where $c_{\text{start}}$ is the initial configuration and $c_{\text{accept}}$ is the accepting configuration. This runs in $\mathsf{DSPACE}(s(n)^2)$.
\end{proof}

\begin{corollary}
$\mathsf{PSPACE} = \mathsf{NPSPACE}$. (Unlike time, non-determinism gives at most quadratic space advantage.)
\end{corollary}

\begin{corollary}
$\mathsf{NL} \subseteq \mathsf{DSPACE}(\log^2 n)$. In particular, $\mathsf{NL} \subseteq \mathsf{P}$.
\end{corollary}

\subsection{Relationships Between Classes}
\label{sec:complexity:relationships}

\begin{theorem}
$\mathsf{L} \subseteq \mathsf{NL} \subseteq \mathsf{P} \subseteq \mathsf{NP} \subseteq \mathsf{PSPACE} \subseteq \mathsf{EXP}$.
\end{theorem}

\begin{proof}[Key Inclusions]
\begin{itemize}
\item $\mathsf{NL} \subseteq \mathsf{P}$: Configuration graph of logspace NTM has poly size; reachability in P.
\item $\mathsf{P} \subseteq \mathsf{NP}$: Deterministic is special case of non-deterministic.
\item $\mathsf{NP} \subseteq \mathsf{PSPACE}$: Try all non-deterministic branches in polynomial space (DFS).
\item $\mathsf{PSPACE} \subseteq \mathsf{EXP}$: At most $2^{O(s(n))}$ configurations; simulate in exp time.
\end{itemize}
\end{proof}

\begin{remark}
It is known that $\mathsf{L} \neq \mathsf{PSPACE}$ (by space hierarchy), $\mathsf{P} \neq \mathsf{EXP}$ (by time hierarchy), and $\mathsf{NL} \neq \mathsf{PSPACE}$ (by space hierarchy applied to NSPACE). But we do not know any of: $\mathsf{L} = \mathsf{P}$, $\mathsf{P} = \mathsf{NP}$, $\mathsf{NP} = \mathsf{PSPACE}$, $\mathsf{PSPACE} = \mathsf{EXP}$. These are the central open questions.
\end{remark}

\begin{theorem}[Padding Argument]
If $\mathsf{P} = \mathsf{NP}$, then $\mathsf{EXP} = \mathsf{NEXP}$. More generally, for any ``reasonable'' time classes, collapsing smaller classes collapses larger ones.
\end{theorem}

\begin{proof}[Sketch]
Given a language $L \in \mathsf{NEXP}$ decided by a non-deterministic TM in $2^{n^k}$ time, define its \emph{padded version} $L' = \{ x10^{2^{|x|^k}} \mid x \in L \}$. $L'$ is in $\mathsf{NP}$ (the padding provides enough time to simulate the exponential computation). If $\mathsf{P} = \mathsf{NP}$, then $L' \in \mathsf{P}$, and therefore $L \in \mathsf{EXP}$ by reversing the padding. So $\mathsf{EXP} = \mathsf{NEXP}$.
\end{proof}

\begin{theorem}[Ladner's Theorem (1975)]
If $\mathsf{P} \neq \mathsf{NP}$, then there exist problems in $\mathsf{NP} \setminus \mathsf{P}$ that are \emph{not} NP-complete (NP-intermediate).
\end{theorem}

\begin{proof}[Proof Sketch]
The proof constructs a language $L$ by a stage-wise diagonalization. Let $f(n) = n^{\log n}$ (or any function that grows faster than any polynomial but slower than exponentials). Define an NP problem:
\[
\text{SAT-HALT} = \{ \phi 10^{f(|\phi|)} \mid \phi \in \text{SAT} \}
\]
where the padding $10^{f(|\phi|)}$ ensures the input length is stretched. 

We need to ensure SAT-HALT is:
\begin{itemize}
\item In NP: The verifier strips the padding and checks $\phi$ is satisfiable.
\item Not in P: If it were in P, we could solve SAT in $O(p(f(n)))$ time — but since $f$ grows super-polynomially, this would put SAT in P, contradicting $\mathsf{P} \neq \mathsf{NP}$.
\item Not NP-complete: Any polynomial-time reduction from SAT to SAT-HALT would have to deal with the padding — the output size must be polynomial in the input. The super-polynomial padding prevents this, because the reduction cannot ``stretch'' the input enough to simulate the padding requirement.
\end{itemize}

The actual construction uses a delicate \emph{delayed diagonalization}: we define a language that alternates between behaving like SAT and like a polynomial-time language, with the switch points determined by a function $g(n)$ that grows slowly enough to keep the problem out of P but fast enough to prevent NP-completeness.
\end{proof}

\begin{remark}[NP-Intermediate Candidates]
If $\mathsf{P} \neq \mathsf{NP}$, then Ladner's theorem guarantees a rich hierarchy within NP. Leading candidates for NP-intermediate status include: Graph Isomorphism, Factoring, the Shortest Vector Problem (SVP) in lattices, and the Minimum Circuit Size Problem (MCSP). None is known to be in P or NP-complete.
\end{remark}

\begin{theorem}[Relativization (Baker-Gill-Solovay 1975)]
There exist oracles $A$ and $B$ such that $\mathsf{P}^A = \mathsf{NP}^A$ and $\mathsf{P}^B \neq \mathsf{NP}^B$. Therefore, any proof resolving P vs NP must be \emph{non-relativizing}.
\end{theorem}

\begin{proof}[Sketch]
For $A$, take a PSPACE-complete language. Then $\mathsf{NP}^A \subseteq \mathsf{PSPACE} \subseteq \mathsf{P}^A$ since $\mathsf{P}^{\mathsf{PSPACE}} = \mathsf{PSPACE}$. For $B$, define an oracle that encodes a unary language $L_B = \{ 1^n \mid \text{the } n\text{-th TM doesn't accept } 1^n \text{ in } 2^n \text{ steps with oracle } B \}$. The standard diagonalization shows $\mathsf{P}^B \neq \mathsf{NP}^B$.
\end{proof}

\subsection{The Extended Church-Turing Thesis}
\label{sec:complexity:ect}

\begin{quote}
\textbf{Extended Church-Turing Thesis:} Any ``reasonable'' model of computation can be simulated on a probabilistic Turing machine with at most polynomial overhead.
\end{quote}

This makes P, BPP, etc. \emph{physically} meaningful — not just TM artifacts. Quantum computers challenge this (BQP $\not\subseteq$ BPP?).

The invariance thesis, formalized by Pippenger and Fischer (1979), states that all ``reasonable'' sequential models (RAM machines, pointer machines, $\lambda$-calculus, cellular automata) simulate each other with polynomial overhead. The key is that Turing machines with multiple tapes can simulate $k$-tape TMs with $O(\log k)$ overhead, and RAM machines simulate TMs with $O(n^2)$ overhead. This robustness justifies P as a model-independent class.

\subsection{Circuit Complexity and Non-Uniform Models}
\label{sec:complexity:circuits}

Circuit complexity studies the size/depth of Boolean circuits computing functions. Unlike TMs, circuits are \emph{non-uniform}: a different circuit can be used for each input length.

\begin{definition}[Circuit Classes]
\begin{itemize}
\item $\mathsf{AC}^0$: constant depth, unbounded fan-in AND/OR/NOT gates.
\item $\mathsf{AC}^0[p]$: $\mathsf{AC}^0$ with MOD$_p$ gates.
\item $\mathsf{TC}^0$: constant depth, majority (threshold) gates.
\item $\mathsf{NC}^1$: log depth, bounded fan-in.
\item $\mathsf{NC}$: polylog depth, bounded fan-in.
\item $\mathsf{P}/\mathsf{poly}$: polynomial-size circuit families (non-uniform).
\end{itemize}
\end{definition}

\begin{theorem}[Furst-Saxe-Sipser 1984; Ajtai 1983; Yao 1985; H{\aa}stad 1986]
Parity $\notin \mathsf{AC}^0$. Specifically, any depth-$d$ circuit computing parity on $n$ bits requires size $\exp(\Omega(n^{1/(d-1)}))$.
\end{theorem}

\begin{proof}[Proof Sketch via the Switching Lemma]
The proof uses the \emph{switching lemma}: any depth-2 circuit (a CNF or DNF) can be ``switched'' to the opposite form with high probability after a random restriction (setting most variables randomly). By iterating this lemma, any $\mathsf{AC}^0$ circuit can be reduced to a shallow decision tree that depends on few variables.

\textbf{Step 1:} Randomly restrict variables: for each variable independently, set it to 0 or 1 with probability $p$, leave it free with probability $1-2p$. 

\textbf{Step 2:} The switching lemma guarantees that a CNF of width $w$ becomes a DNF of width $w$ with high probability after restriction. By applying this bottom-up through the circuit, we reduce each layer.

\textbf{Step 3:} After $d-1$ applications, the circuit collapses to a function depending on at most $O(\log n)$ variables — a constant function relative to the number of remaining free variables.

\textbf{Step 4:} But parity of $n$ bits, after a random restriction leaving $m$ free variables, is still parity on $m$ bits (uniformly random). Parity on $m$ variables depends on all $m$ of them. So for a circuit to compute parity, it must have depth at least $\Omega(\log n / \log \log n)$ or exponential size.

The tight lower bound, due to H{\aa}stad: any depth-$d$ circuit for parity has size $\exp(\Omega(n^{1/(d-1)}))$.
\end{proof}

\begin{example}[Why Parity Is Hard for AC$^0$]
Parity is the function $\text{XOR}(x_1, \dots, x_n) = 1$ iff an odd number of inputs are 1. A DNF for parity requires $2^{n-1}$ terms (all odd-parity assignments). A CNF requires $2^{n-1}$ clauses. The switching lemma shows this blowup is unavoidable at any constant depth — small-depth circuits cannot ``aggregate'' information across all $n$ inputs without exponential growth.
\end{example}

\begin{theorem}[Razborov-Smolensky 1987]
MOD$_p$ gates cannot compute MOD$_q$ for distinct primes $p, q$. $\mathsf{AC}^0[p]$ lower bounds.
\end{theorem}

\begin{proof}[Sketch]
The proof uses polynomial approximations over finite fields. Suppose a circuit with MOD$_p$ gates computes MOD$_q$. Each MOD$_p$ gate can be approximated by a low-degree polynomial over $\mathbb{F}_p$. Composing these approximations gives a low-degree polynomial approximating MOD$_q$. But MOD$_q$ requires degree $\Omega(n)$ to approximate over $\mathbb{F}_p$ (because MOD$_q$ is not a polynomial function over $\mathbb{F}_p$ when $p \neq q$). This gives a contradiction for sufficiently large $n$.
\end{proof}

\begin{theorem}[Adleman 1978]
$\mathsf{BPP} \subseteq \mathsf{P}/\mathsf{poly}$. Any randomized algorithm can be derandomized with polynomial-size non-uniform circuits.
\end{theorem}

\begin{proof}[Sketch]
Let $L \in \mathsf{BPP}$ with error $\le 1/3$. By running the algorithm $O(n)$ times and taking the majority vote, we reduce error to $< 2^{-n}$. By the union bound, there exists a single sequence of random bits that works for all $2^n$ inputs of length $n$. Hard-wire this sequence as advice. The resulting circuit is polynomial-size.
\end{proof}

\begin{theorem}[Natural Proofs Barrier (Razborov-Rudich 1997)]
If strong pseudorandom generators exist (i.e., one-way functions exist), then no ``natural'' property can prove superpolynomial circuit lower bounds for NP. A property $\mathcal{P}$ is \emph{natural} if:
\begin{itemize}
\item \textbf{Constructive:} $\mathcal{P}$ is decidable in $2^{O(n)}$ time (can be checked by a circuit of size $2^{O(n)}$).
\item \textbf{Large:} $\mathcal{P}$ holds for at least a $1/\mathsf{poly}(n)$ fraction of all $n$-input Boolean functions.
\end{itemize}
A natural property \emph{useful against} $\mathsf{P}/\mathsf{poly}$ is one that holds for all functions in $\mathsf{P}/\mathsf{poly}$ and fails for some function (which would separate NP from $\mathsf{P}/\mathsf{poly}$).
\end{theorem}

\begin{proof}[Explanation]
Suppose there exists a natural property $\mathcal{P}$ useful against $\mathsf{P}/\mathsf{poly}$. Being constructive, $\mathcal{P}$ can be computed by a circuit of size $2^{O(n)}$. Being large, a random function satisfies $\mathcal{P}$ with probability $1/\mathsf{poly}(n)$. 

Now consider a candidate pseudorandom generator $G: \{0,1\}^{n^c} \to \{0,1\}^{2^n}$, computable in $\mathsf{DTIME}(2^{O(n)})$, that fools circuits of size $2^{O(n)}$. Since $\mathcal{P}$ is constructible, $G$ must fool it: $\Pr_{x}[ \mathcal{P}(G(x)) ] \approx \Pr_{y \in \{0,1\}^{2^n}}[ \mathcal{P}(y) ] \ge 1/\mathsf{poly}(n)$. 

But the image of $G$ has size $2^{n^c}$, which is $\mathsf{P}/\mathsf{poly}$ (each output corresponds to a polynomial-size circuit by the PRG construction). If $\mathcal{P}$ holds for all $\mathsf{P}/\mathsf{poly}$ functions, then $\mathcal{P}(G(x))$ holds with probability 1. This contradicts the approximation if $1/\mathsf{poly}(n) < 1$. Formalizing this yields: if one-way functions exist, natural proofs cannot separate $\mathsf{NP}$ from $\mathsf{P}/\mathsf{poly}$.

This barrier explains why we cannot prove P $\neq$ NP via circuit lower bounds — unless we find non-constructive or non-large properties. Most known circuit lower bound techniques (including the switching lemma and Razborov-Smolensky) give natural properties, which means they cannot separate P from NP.
\end{proof}

The Razborov-Rudich barrier is not just a technical result — it is a deep insight about the limits of current mathematical reasoning. Most lower bound techniques (random restrictions, polynomial approximation, combinatorial discrepancy) produce ``natural'' properties. The barrier says: if cryptography is possible (one-way functions exist), then these techniques \emph{cannot} prove NP $\not\subseteq$ P/poly. To separate P from NP, we need entirely new methods — perhaps those that exploit the \emph{uniformity} of computation (the fact that a single algorithm works for all input sizes). This is why progress on P vs NP has been so slow: the obvious tools are provably inadequate.

\subsection{Derandomization and Pseudorandomness}
\label{sec:complexity:derand}

Derandomization asks: can every randomized algorithm be made deterministic without losing efficiency? The answer, surprisingly, depends on circuit lower bounds.

\begin{definition}[Pseudorandom Generator (PRG)]
A function $G: \{0,1\}^s \to \{0,1\}^n$ is a \emph{pseudorandom generator} against a class of circuits $\mathcal{C}$ if for every $C \in \mathcal{C}$:
\[
\left| \Pr_{x \in \{0,1\}^s}[C(G(x)) = 1] - \Pr_{y \in \{0,1\}^n}[C(y) = 1] \right| \le \epsilon
\]
The PRG \emph{fools} $\mathcal{C}$: no circuit in $\mathcal{C}$ can distinguish the output of $G$ from truly random bits.
\end{definition}

\begin{theorem}[Nisan-Wigderson 1994]
Let $f: \{0,1\}^n \to \{0,1\}$ be a function hard for circuits of size $2^{\epsilon n}$ for some $\epsilon > 0$. Then there exists a pseudorandom generator $G_{\text{NW}}: \{0,1\}^{O(\log n)} \to \{0,1\}^n$ that fools circuits of size $n$.
\end{theorem}

\begin{proof}[Sketch of NW Construction]
The NW generator uses a \emph{hardness-amplified} function $f$ and an \emph{(n, k)-design} — a collection of subsets $S_1, \dots, S_n \subseteq [\ell]$ (where $\ell = O(\log n)$) such that each $|S_i| = m$ and $|S_i \cap S_j|$ is small. The generator is:
\[
G_{\text{NW}}(x) = f(x|_{S_1}) \circ f(x|_{S_2}) \circ \cdots \circ f(x|_{S_n})
\]
where $x \in \{0,1\}^\ell$ is the seed and $x|_S$ denotes the bits of $x$ indexed by $S$.

\textbf{Proof by contradiction:} Suppose a circuit $C$ of size $n$ distinguishes $G_{\text{NW}}(x)$ from uniform. By a hybrid argument, there exists a position $i$ and a circuit that predicts the $i$-th bit given the first $i-1$ bits with advantage $> \epsilon/n$. Using the design property (small intersections), this yields a circuit that computes $f$ on a fresh input with nontrivial advantage, contradicting the hardness of $f$.

The existence of such $f$ (exponential circuit complexity, computable in $\mathsf{E}$) implies $\mathsf{P} = \mathsf{BPP}$ via a universal derandomization.
\end{proof}

\begin{theorem}[Impagliazzo-Wigderson 1997]
If $\mathsf{E}$ requires exponential-size circuits (i.e., $\mathsf{E} \not\subseteq \mathsf{SIZE}(2^{\epsilon n})$), then $\mathsf{P} = \mathsf{BPP}$.
\end{theorem}

\begin{proof}[High-Level Structure]
\begin{enumerate}
\item Start with a problem in $\mathsf{E}$ that requires circuits of size $2^{\epsilon n}$. By a hardness amplification lemma (Yao's XOR lemma), we can construct a function $f$ that is $2^{\epsilon n}$-hard for circuits.
\item Use the NW construction to build a PRG $G: \{0,1\}^{O(\log n)} \to \{0,1\}^n$ that fools circuits of size $n$ with seed length $O(\log n)$.
\item To derandomize a BPP algorithm $A(x, r)$ (using $n$ random bits), enumerate all $2^{c \log n} = n^c$ seeds and check if \emph{majority} of $A(x, G(s))$ accept. This runs in polynomial time.
\end{enumerate}
Thus every BPP algorithm can be simulated in P, assuming sufficiently strong circuit lower bounds.
\end{proof}

The IW theorem establishes a deep duality: hardness (worst-case circuit lower bounds) is equivalent to randomness (the ability to derandomize). This is one of the most beautiful connections in complexity theory. It says: either some problem is truly hard (exponential circuit complexity), or randomness gives no power at all (BPP = P). We cannot have both: strong pseudorandomness \emph{requires} and \emph{implies} hardness.

\subsection{Quantum Complexity}
\label{sec:complexity:quantum}

Quantum computation challenges the Extended Church-Turing Thesis by proposing a physically realizable model that may offer exponential speedups.

\begin{definition}[BQP]
$\mathsf{BQP}$ = languages decidable by a uniform family of quantum circuits in polynomial time with bounded error ($\le 1/3$). The circuit uses polynomial-size quantum gates (Hadamard, CNOT, phase gates) and ends with a measurement.
\end{definition}

\begin{definition}[Key Quantum Gates]
\begin{itemize}
\item \textbf{Hadamard:} $H|0\rangle = \frac{1}{\sqrt{2}}(|0\rangle + |1\rangle)$, $H|1\rangle = \frac{1}{\sqrt{2}}(|0\rangle - |1\rangle)$. Creates superposition.
\item \textbf{CNOT:} Controlled-NOT flips the target qubit if the control qubit is 1.
\item \textbf{Phase:} $R_\theta|0\rangle = |0\rangle$, $R_\theta|1\rangle = e^{i\theta}|1\rangle$. Introduces relative phase.
\end{itemize}
\end{definition}

\begin{theorem}[Shor 1994]
Factoring and Discrete Log are in $\mathsf{BQP}$.
\end{theorem}

\begin{proof}[Sketch of Shor's Algorithm]
Shor's algorithm reduces factoring to order-finding, then solves order-finding using the Quantum Fourier Transform (QFT).

\textbf{Step 1 — Reduction:} To factor $N$, pick random $a < N$. Compute $g = \gcd(a, N)$. If $g > 1$, we found a factor. Otherwise, find the order $r$ of $a$ modulo $N$ — the smallest $r$ such that $a^r \equiv 1 \pmod{N}$.

\textbf{Step 2 — Quantum order-finding:} Use two quantum registers. Initialize to uniform superposition over $q$ basis states ($q = 2^m$, $N^2 \le q < 2N^2$):
\[
\frac{1}{\sqrt{q}} \sum_{x=0}^{q-1} |x\rangle |0\rangle
\]

\textbf{Step 3 — Modular exponentiation:} Compute $a^x \bmod N$ in the second register:
\[
\frac{1}{\sqrt{q}} \sum_{x=0}^{q-1} |x\rangle |a^x \bmod N\rangle
\]

\textbf{Step 4 — QFT:} Apply the Quantum Fourier Transform to the first register:
\[
\frac{1}{q} \sum_{x=0}^{q-1} \sum_{y=0}^{q-1} e^{2\pi i xy/q} |y\rangle |a^x \bmod N\rangle
\]

\textbf{Step 5 — Measurement:} Measuring the first register gives $y$ with probability concentrated near values where $y/q \approx k/r$ for some integer $k$. Use continued fractions to extract $r$ from $y/q$.

\textbf{Step 6 — Factorization:} Once $r$ is known, with high probability $\gcd(a^{r/2} \pm 1, N)$ yields a non-trivial factor of $N$.
\end{proof}

\begin{theorem}[Grover 1996]
Search in an unstructured database of size $N$ requires $\Omega(\sqrt{N})$ quantum queries. The optimal quantum search algorithm achieves this bound.
\end{theorem}

\begin{proof}[Sketch of Grover's Algorithm]
Grover's algorithm uses amplitude amplification. Start with uniform superposition $|\psi\rangle = \frac{1}{\sqrt{N}} \sum_{x=0}^{N-1} |x\rangle$. Define two operations: (1) the oracle $O_f$ that flips the phase of the target state, and (2) the Grover diffusion operator $D = 2|\psi\rangle\langle\psi| - I$, which reflects about $|\psi\rangle$. Applying $O_f D$ repeatedly $O(\sqrt{N})$ times rotates the state toward the target. Measurement then gives the target with high probability. The lower bound $\Omega(\sqrt{N})$ follows from the fact that quantum operations are unitary and information from the oracle is limited by the angle between initial and target states.
\end{proof}

\begin{theorem}[Bennett-Bernstein-Brassard-Vazirani 1997]
$\mathsf{BQP} \subseteq \mathsf{PSPACE}$.
\end{theorem}

\begin{proof}[Sketch]
The acceptance probability of a quantum circuit can be expressed as a sum of exponentials of matrix entries. By summing over all computational paths (which are exponentially many but can be enumerated in polynomial space using depth-first search), we can compute the probability exactly. The key: quantum amplitudes are complex numbers, but deciding whether a given outcome has nonzero probability reduces to evaluating a polynomial-sized sum-of-products, which is in PSPACE.
\end{proof}

Open questions: $\mathsf{BQP}$ vs $\mathsf{PH}$? $\mathsf{BQP}$ vs $\mathsf{NP}$? Is factoring in P? Does quantum supremacy (random circuit sampling) imply $\mathsf{BQP} \not\subseteq \mathsf{PH}$?

Google's 2019 experiment (Sycamore processor) claimed to demonstrate that random circuit sampling is beyond classical simulation. The complexity-theoretic basis: if sampling from the output distribution of random quantum circuits were in $\mathsf{BPP}$, then $\mathsf{PH}$ would collapse to the third level. This is a conditional hardness result: being in BPP is unlikely because it would have implausible consequences for the polynomial hierarchy. Thus, quantum supremacy experiments are complexity-theoretic, not just engineering feats.

\subsection{Fine-Grained Complexity}
\label{sec:complexity:fine}

Instead of ``P vs NP,'' asks: ``Is this $O(n^2)$ or $O(n^{3-\epsilon})$?'' While P vs NP separates polynomial from exponential, fine-grained complexity distinguishes among polynomial runtimes.

\begin{definition}[Strong Exponential Time Hypothesis (SETH)]
For every $\epsilon > 0$, there exists $k$ such that $k$-SAT cannot be solved in $O(2^{(1-\epsilon)n})$ time. Equivalently, SAT requires essentially $2^n$ time — no subexponential algorithm exists even with unbounded clause length.
\end{definition}

\begin{definition}[Orthogonal Vectors Hypothesis (OVH)]
Given two sets of $n$ vectors in $\{0,1\}^d$, deciding if there's an orthogonal pair requires $n^{2-o(1)}$ time (for $d = \omega(\log n)$).
\end{definition}

\begin{theorem}[Equivalences under OVH and SETH]
Under SETH, many problems require quadratic time:
\begin{itemize}
\item Edit Distance, LCS (Longest Common Subsequence)
\item Fr\'echet Distance
\item Dynamic Time Warping
\item Pattern Matching with Wildcards
\end{itemize}
Each of these problems has an $O(n^2)$ dynamic programming algorithm that is believed to be optimal.
\end{theorem}

\begin{example}[Edit Distance Lower Bound]
Computing the edit distance (Levenshtein distance) between two strings of length $n$ can be done in $O(n^2)$ time via DP. A sub-quadratic algorithm ($O(n^{2-\epsilon})$) would refute SETH. This is shown by encoding a SAT instance as a pair of strings such that edit distance encodes the satisfiability of the instance. The reduction is delicate: each clause becomes a gadget in the string, and the edit distance calculation simulates the search for a satisfying assignment.
\end{example}

\begin{definition}[3SUM Hypothesis]
3SUM (given $n$ integers, find three summing to 0) requires $n^{2-o(1)}$ time.
\end{definition}

\begin{definition}[APSP Equivalence]
All-Pairs Shortest Paths (APSP) is equivalent to (min,+) matrix multiplication, negative triangle detection, and several other graph problems. Any $O(n^{3-\epsilon})$ algorithm for one yields such an algorithm for all.
\end{definition}

\begin{theorem}[Williams 2014]
If $\mathsf{NEXP} \not\subseteq \mathsf{P}/\mathsf{poly}$, then some problem in $\mathsf{NTIME}(2^{O(n)})$ requires $n^{\omega(\log n)}$ ACC$^0$ circuit size — a non-uniform fine-grained lower bound. This connects non-uniform lower bounds (circuit complexity) to uniform separations (NEXP vs P/poly).
\end{theorem}

\subsection{Parameterized Complexity}
\label{sec:complexity:param}

Parameterized complexity studies problems with a natural parameter $k$ (e.g., solution size, treewidth, dimension). The goal: is the problem tractable when $k$ is small, even if $n$ is large?

\begin{definition}[Fixed-Parameter Tractable (FPT)]
A problem is FPT if it can be solved in $f(k) \cdot n^{O(1)}$ time, where $k$ is a parameter (e.g., solution size, treewidth). The function $f$ may grow super-polynomially (e.g., $2^k$), but the exponent of $n$ is independent of $k$.
\end{definition}

\begin{definition}[Treewidth]
A graph has treewidth $t$ if it can be decomposed into a tree of overlapping bags, each containing at most $t+1$ vertices. Treewidth measures how ``tree-like'' a graph is. Many NP-hard problems become FPT parameterized by treewidth via dynamic programming on tree decompositions.
\end{definition}

\begin{theorem}[Downey-Fellows 1999]
$k$-Vertex Cover is FPT ($O(1.2738^k + kn)$, Chen-Kanj-Xia 2010). $k$-Clique is W[1]-complete (not FPT unless FPT = W[1]).
\end{theorem}

\begin{proof}[Sketch: Vertex Cover]
The classic FPT algorithm uses bounded search tree branching. Pick any edge $(u,v)$. Any vertex cover must include at least one of $u$ or $v$. Branch on both possibilities, removing covered edges, and recurse with $k-1$. This yields $O(2^k)$ branching. The improved algorithm uses more careful case analysis (crown reduction, kernelization) to achieve $O(1.2738^k)$.

\textbf{Kernelization:} First reduce the instance to a smaller ``kernel'' of size $O(k^2)$ using the following rules: (1) remove isolated vertices; (2) if a vertex has degree $> k$, it must be in the cover (add it, decrease $k$); (3) otherwise the remaining graph has at most $k^2 + k$ edges. Then apply the branching on the kernel.
\end{proof}

The W-hierarchy: $\mathsf{FPT} \subseteq \mathsf{W}[1] \subseteq \mathsf{W}[2] \subseteq \dots \subseteq \mathsf{XP}$.

\begin{definition}[W-hierarchy]
$\mathsf{W}[t]$ is defined via the weft of a circuit — roughly, the number of alternations between large fan-in AND/OR gates on any path from input to output. Parameterized reductions preserve the parameter up to a polynomial. Key complete problems:
\begin{itemize}
\item $\mathsf{W}[1]$-complete: $k$-Clique, $k$-Independent Set
\item $\mathsf{W}[2]$-complete: $k$-Dominating Set, $k$-Set Cover
\end{itemize}
If any $\mathsf{W}[t]$-complete problem is in FPT, then $\mathsf{FPT} = \mathsf{W}[t]$, which would imply $\mathsf{FPT} = \mathsf{W}[t]$ for all $t$. This is believed false.
\end{definition}

\begin{example}[Parameterized Intractability: $k$-Clique]
$k$-Clique is complete for $\mathsf{W}[1]$ via a parameterized reduction from Independent Set. The brute-force $O(n^k)$ algorithm is not FPT because $k$ appears in the exponent. Showing $k$-Clique $\in$ FPT would imply $\mathsf{FPT} = \mathsf{W}[1]$, which is as unlikely as P = NP. The parameterized lens reveals a fine-grained hierarchy within NP-hard problems: some admit $f(k) \cdot n^{O(1)}$ algorithms (Vertex Cover), while others inherently require $n^{\Omega(k)}$ (Clique).
\end{example}

\subsection{Proof Complexity}
\label{sec:complexity:proof}

Proof complexity studies the length of proofs in propositional proof systems. It connects complexity theory to logic: lower bounds on proof length imply computational lower bounds.

\begin{definition}[Propositional Proof Systems]
\begin{itemize}
\item \textbf{Resolution:} Derive the empty clause from a CNF formula using the rule $(C \lor x), (D \lor \lnot x) \vdash (C \lor D)$. Refutation = deriving the empty clause (contradiction).
\item \textbf{Cutting Planes:} Linear inequalities over integers. Derive $0 \ge 1$ from a set of 0-1 inequalities representing clauses. Uses addition and division rules.
\item \textbf{Polynomial Calculus:} Polynomial equations over a field. Derive $1 = 0$ from $x_i^2 - x_i = 0$ and clause polynomials. Uses algebraic reasoning.
\item \textbf{Sum-of-Squares (SoS):} Polynomial inequalities with sum-of-squares certificates. Powerful for optimization (approximation algorithms).
\item \textbf{Extended Frege:} Frege with definitions — allows introducing new variables for subformulas. The most powerful natural system.
\end{itemize}
\end{definition}

\begin{theorem}[Cook-Reckhow 1979]
$\mathsf{NP} = \mathsf{coNP}$ iff there exists a polynomially-bounded proof system (all tautologies have poly-size proofs).
\end{theorem}

\begin{proof}[Sketch]
($\Rightarrow$) If $\mathsf{NP} = \mathsf{coNP}$, then $\mathsf{TAUT}$ (the set of tautologies) is in $\mathsf{NP}$. The NP verifier for $\mathsf{TAUT}$ acts as a polynomially-bounded proof system: the certificate is a proof that the formula is a tautology.

($\Leftarrow$) If a polynomially-bounded proof system exists, then $\mathsf{TAUT} \in \mathsf{NP}$: on input $\phi$, guess a proof and verify it. Since $\mathsf{TAUT}$ is coNP-complete, $\mathsf{coNP} \subseteq \mathsf{NP}$ (and symmetrically $\mathsf{NP} \subseteq \mathsf{coNP}$), so $\mathsf{NP} = \mathsf{coNP}$.
\end{proof}

\begin{theorem}[Resolution Width Lower Bound (Ben-Sasson-Wigderson 2001)]
If a CNF formula $F$ has resolution complexity $S$, then there exists a proof of width $\le O(\sqrt{n \log S})$. Conversely, any formula requiring width $w$ requires proof size $\exp(\Omega(w^2/n))$. Thus width lower bounds imply size lower bounds.
\end{theorem}

\begin{proof}[Sketch]
The upper bound: any resolution proof can be ``compressed'' to a proof with bounded width by a careful extraction of a subproof. The lower bound: define a graph whose vertices are clauses; clause width corresponds to the number of variables. The size-width relationship shows that exponential-size lower bounds follow from linear-width lower bounds. This is used to prove that the pigeonhole principle requires exponential resolution proofs ($\mathsf{PHP}_n^{n+1}$ has no poly-size resolution refutation).
\end{proof}

\begin{example}[Pigeonhole Principle: Proof Complexity Lower Bound]
The pigeonhole principle ($n+1$ pigeons in $n$ holes forces some hole to contain 2 pigeons) is a tautology. Encoding it as a CNF yields a formula $\mathsf{PHP}_n^{n+1}$ with $O(n^2)$ clauses. Resolution requires $\exp(\Omega(n))$ steps to refute it, despite the formula being small. This shows resolution cannot efficiently prove even simple counting principles. In contrast, Extended Frege can prove $\mathsf{PHP}$ in polynomial size using definitions that introduce counting predicates.
\end{example}

\begin{remark}[Proof Complexity and Circuit Lower Bounds]
Feasible interpolation: if a proof system has feasible interpolation, then lower bounds for the proof system imply circuit lower bounds. This connects proof complexity (non-uniform) to circuit complexity. For example, Cutting Planes has feasible interpolation, which allowed the proof that $\mathsf{NP} \neq \mathsf{coNP}$ cannot be proved in Cutting Planes (unless certain cryptographic assumptions fail). Resolution also has feasible interpolation (via monotone circuit lower bounds). This is why proof complexity lower bounds are hard: they would imply circuit lower bounds, which the Natural Proofs barrier suggests are beyond current techniques.
\end{remark}

%======================================================================
% SECTION 6: CONSEQUENCES
%======================================================================
\section{Consequences: What Became Possible}
\label{sec:complexity:consequences}

\begin{enumerate}
\item \textbf{Algorithm Design as Optimization:} Now we ask: ``Is this problem in P? If not, can we approximate it? Is it fixed-parameter tractable? Can we solve it for small $k$ quickly?'' The complexity class \emph{guides} the algorithmic approach — choosing approximation, heuristics, parameterization, or exact exponential algorithms.

\item \textbf{Cryptography (Chapter 10):} Modern crypto \emph{requires} P $\neq$ NP (actually, average-case hardness, one-way functions). If P = NP, most crypto collapses. The complexity-theoretic foundation of cryptography is the existence of problems that are hard on average, not just in the worst case. This is why cryptography is not just ``applied complexity'' — it introduces fundamentally new requirements (average-case hardness, one-way functions).

\item \textbf{Optimization:} NP-completeness explains why TSP, SAT, Knapsack, Graph Coloring resist exact polynomial algorithms. We use: approximation algorithms (with provable guarantees), ILP solvers (branch-and-cut), SAT solvers (CDCL), local search (no guarantees but practical), parameterized algorithms (exponential in $k$ only), and exact exponential algorithms ($O(c^n)$ with small $c$).

\item \textbf{Proof Complexity:} Cook-Reckhow (1979): propositional proof systems. Lower bounds on proof length $\leftrightarrow$ lower bounds on computation. P vs NP $\leftrightarrow$ existence of polynomially-bounded proof systems. Automated theorem provers (SAT solvers) are essentially searching for Resolution proofs — their power and limitations are exactly those of the proof system.

\item \textbf{Descriptive Complexity (Fagin 1974):} $\mathsf{NP} = \exists$SO (existential second-order logic). $\mathsf{P} = \mathsf{FO}$(LFP) (fixed-point logic). Complexity = logic expressiveness. This means P vs NP is equivalent to: can every property expressible in second-order logic also be expressed in first-order logic with fixed-point operators? This reformulation connects complexity to model theory.

\item \textbf{Quantum Complexity (BQP):} Shor's algorithm (1994): factoring in BQP. Not known to be in P. BQP vs PH is a major open question. If BQP $\not\subseteq$ PH, then quantum computers can solve problems beyond classical heuristic reasoning.

\item \textbf{Derandomization:} $\mathsf{P} = \mathsf{BPP}$? (Impagliazzo-Wigderson 1997: if $\mathsf{E}$ requires exponential circuits, then $\mathsf{P} = \mathsf{BPP}$.) Pseudorandom generators connect hardness to randomness. This has practical implications: we can often replace random bits with the output of a cryptographic PRG and be confident of correctness.

\item \textbf{Learning Theory:} PAC learning (Valiant 1984) is intimately connected to circuit complexity. A concept class is efficiently PAC-learnable iff there is a compression scheme for it. The equivalence: learnability $\iff$ Occam's Razor. Cryptographic lower bounds for learning (under worst-case hardness assumptions) show that some concept classes (like DNF formulas) are not efficiently learnable.

\end{enumerate}

%======================================================================
% SECTION 7: PHILOSOPHICAL SIGNIFICANCE
%======================================================================
\section{Philosophical Significance: What It Revealed About Limits of Formal Systems}
\label{sec:complexity:philosophy}

\begin{enumerate}
\item \textbf{Difficulty Is Intrinsic, Not Algorithmic:} A problem's complexity is a property of the \emph{problem}, not the \emph{algorithm}. This is a Platonic view: the ``hardness'' of TSP exists independently of humans. Like the primes, the complexity classes were there waiting to be discovered — they existed as mathematical structures long before we named them.

\item \textbf{Verification vs. Generation:} NP = ``easy to check, hard to find.'' This asymmetry (P vs NP) is a fundamental fact about \emph{solutions} vs. \emph{proofs}. Creativity (finding) is harder than criticism (verifying) — a mathematical truth. If P = NP, then creativity and verification are equivalent: every puzzle whose solution is easy to check is easy to solve. This would have profound implications for mathematics (automated theorem proving), science (hypothesis generation), and artificial intelligence.

\item \textbf{The Map of the Computable Universe:} Complexity theory is the \emph{geography} of computation. P, NP, PSPACE, EXP are continents. Reductions are roads. Complete problems are capitals. We are mapping a landscape that exists whether we explore it or not. The completeness phenomenon — that seemingly unrelated problems (SAT, TSP, 3-Coloring, Subset Sum) are the same problem in disguise — reveals a deep unity of computational difficulty.

\item \textbf{Feasibility Has a Phase Transition:} Polynomial vs. exponential is not a gradient — it's a \emph{phase boundary}. Problems just above the boundary (e.g., $n^{\log n}$) are ``moderately exponential'' — often practically solvable for moderate $n$. The boundary is sharp but the \emph{practical} boundary is fuzzy. Similarly, SAT instances near the clause/variable ratio $4.26$ undergo a phase transition from almost-satisfiable to almost-unsatisfiable.

\item \textbf{Mathematics as Computation:} Proof verification is in P (check each step). Proof \emph{finding} is in NP (guess proof, verify). G\"odel's speed-up theorem: some theorems have exponentially longer proofs in weak systems. Complexity adds a \emph{quantitative} dimension to logic. The P vs NP question is thus about the nature of mathematical creativity itself: is finding a proof no harder than checking one?

\item \textbf{Barriers as Discoveries:} The three great barriers — Relativization (Baker-Gill-Solovay), Natural Proofs (Razborov-Rudich), Algebrization (Aaronson-Wigderson) — teach us about the limits of mathematical reasoning. Each barrier is a theorem that certain proof techniques \emph{cannot} resolve P vs NP. This is meta-mathematics: not about computation, but about what we can prove about computation. The barriers are some of the deepest results in the field.

\end{enumerate}

%======================================================================
% SECTION 8: MODERN LEGACY
%======================================================================
\section{Modern Legacy: Where This Idea Appears Today}
\label{sec:complexity:legacy}

\begin{itemize}
\item \textbf{SAT/SMT Solvers:} CDCL (Conflict-Driven Clause Learning) solves industrial SAT instances with millions of variables. NP-complete in theory; often easy in practice. The ``SAT revolution'' (2000s). Modern solvers combine clause learning, restarts, VSIDS branching heuristic, and preprocessing — techniques that exploit the structure of real-world instances.

\item \textbf{Fine-Grained Complexity:} Not just ``P vs NP'' but ``$O(n^2)$ vs $O(n^{3-\epsilon})$''. Orthogonal Vectors Hypothesis, 3SUM, APSP equivalences. Tight bounds for polynomial-time problems. This is the ``P vs NP for polynomial-time problems'' — understanding the exact exponent of polynomial runtimes.

\item \textbf{Parameterized Complexity:} $k$-Vertex Cover in $O(1.2738^k + kn)$. Downey-Fellows (1999). FPT = Fixed-Parameter Tractable. W-hierarchy for parameterized hardness. The kernelization dichotomy: problems either have polynomial-size kernels (Vertex Cover) or no polynomial kernel unless $\mathsf{NP} \subseteq \mathsf{P/poly}$ (K-CNF, Set Cover).

\item \textbf{Circuit Complexity:} Non-uniform models. $\mathsf{P}/\mathsf{poly}$, $\mathsf{NC}$, $\mathsf{AC}^0$, $\mathsf{TC}^0$. Razborov-Smolensky (1987): MOD$_p$ gates. Natural Proofs barrier (Razborov-Rudich 1997). Recent breakthroughs: lower bounds for $\mathsf{NC}^1$ circuits computing the permanent (but still far from separating P from NP).

\item \textbf{Quantum Supremacy:} Random circuit sampling (Google 2019). BQP vs PH. Complexity theory guides quantum algorithm design and benchmarks quantum hardware. The complexity-theoretic argument for quantum supremacy — that sampling from certain distributions is classically hard — provides evidence that quantum computation is real.

\item \textbf{Meta-Complexity:} Complexity of complexity: MCSP (Minimum Circuit Size Problem), MKTP (Minimum Kt-complexity Problem). Is MCSP NP-complete? (Open since 1970s.) Recent results show that worst-case to average-case reductions can be based on meta-complexity, connecting circuit complexity, cryptography, and meta-complexity in surprising ways.

\item \textbf{Proof Complexity and SAT:} Resolution width/space lower bounds. Feasible interpolation. Connections to circuit lower bounds. Practical SAT solvers correspond to restricted Resolution (tree-like, regular, general). Understanding the limits of SAT solving is understanding the limits of Resolution.

\item \textbf{PAC Learning (Valiant 1984):} Learnability corresponds to compression. Cryptographic hardness implies unlearnability. The connection between complexity and learning: if one-way functions exist, some concept classes (like AC$^0$ under certain distributions) are not efficiently PAC-learnable.

\end{itemize}

\begin{itemize}
\item Chapter 9 (NP-Completeness): The central open problem of this chapter.
\item Chapter 10 (Cryptography): One-way functions $\to$ P $\neq$ NP (necessary but not sufficient).
\item Chapter 11 (Zero Knowledge): ZK $\subseteq$ IP = PSPACE. Complexity classes as proof systems.
\item Chapter 12 (Interactive Proofs): Interactive proofs collapse the polynomial hierarchy; IP = PSPACE.
\item Chapter 13 (Consensus): Distributed complexity (LOCAL, CONGEST models). FLP = async consensus impossible.
\end{itemize}

%======================================================================
% SECTION 9: LESSONS FOR FUTURE SCIENTISTS
%======================================================================
\section{Summary}
\label{sec:complexity:summary}

This chapter developed the theory of computational complexity, classifying problems by their intrinsic resource requirements. The chapter formalized the complexity classes $\mathsf{P}$, $\mathsf{NP}$, $\mathsf{coNP}$, $\mathsf{PSPACE}$, $\mathsf{EXP}$, $\mathsf{NEXP}$, $\mathsf{BPP}$, and $\mathsf{L}$, and established their containment relationships via the hierarchy theorems (Hartmanis--Stearns 1965). Savitch's theorem ($\mathsf{NSPACE}(s(n)) \subseteq \mathsf{DSPACE}(s(n)^2)$) and the Immerman--Szelepcsényi theorem ($\mathsf{NL} = \mathsf{coNL}$) were proven. The chapter discussed the extended Church--Turing thesis and its refutation by quantum algorithms (Shor's factoring algorithm in $\mathsf{BQP}$). Circuit complexity was introduced, with the Furst--Saxe--Sipser and Ajtai results showing $\mathsf{PARITY} \notin \mathsf{AC}^0$, and Håstad's tight switching lemma lower bounds. The polynomial hierarchy $\mathsf{PH}$ was defined via alternating quantifiers, and the relationship $\mathsf{BPP} \subseteq \Sigma_2^{\mathsf{P}} \cap \Pi_2^{\mathsf{P}}$ (Sipser--Gács theorem) was established. The chapter examined derandomization results: the Nisan--Wigderson pseudorandom generator, and the implication that $\mathsf{E} \not\subseteq \mathsf{SIZE}(2^{\epsilon n})$ implies $\mathsf{P} = \mathsf{BPP}$ (Impagliazzo--Wigderson 1997). Fine-grained complexity was discussed via the Orthogonal Vectors Hypothesis and the Strong Exponential Time Hypothesis (SETH). Parameterized complexity was introduced with the W-hierarchy and the concept of fixed-parameter tractability. The natural proofs barrier (Razborov--Rudich 1997) and the relativization barrier (Baker--Gill--Solovay 1975) were analyzed as fundamental obstacles to resolving $\mathsf{P}$ vs $\mathsf{NP}$.

\section*{Key Takeaways}
\begin{itemize}
\item The complexity classes $\mathsf{P}$, $\mathsf{NP}$, $\mathsf{PSPACE}$, $\mathsf{EXP}$, $\mathsf{BPP}$, and $\mathsf{L}$ are related by strict inclusions given by the hierarchy theorems, with $\mathsf{P} \subseteq \mathsf{NP} \subseteq \mathsf{PSPACE} \subseteq \mathsf{EXP}$.
\item Savitch's theorem shows that nondeterministic space can be simulated deterministically with at most a quadratic blowup: $\mathsf{NSPACE}(s) \subseteq \mathsf{DSPACE}(s^2)$.
\item Circuit lower bounds ($\mathsf{AC}^0$ lower bounds by Furst--Saxe--Sipser, Ajtai, Håstad) demonstrate that certain problems require superpolynomial-size circuits.
\item The natural proofs barrier (Razborov--Rudich) and relativization barrier (Baker--Gill--Solovay) show that standard techniques cannot resolve $\mathsf{P}$ vs $\mathsf{NP}$.
\item Open problems include resolving $\mathsf{P}$ vs $\mathsf{NP}$, derandomizing $\mathsf{BPP}$ unconditionally, proving superpolynomial circuit lower bounds for $\mathsf{NP}$, and understanding the relationship between worst-case and average-case complexity (Impagliazzo's five worlds).
\end{itemize}

%======================================================================
% EXERCISES
%======================================================================
% EXERCISES
%======================================================================
\section{Exercises}
\label{sec:complexity:exercises}

\begin{exercise}[Basic]
Prove that if $\mathsf{P} = \mathsf{NP}$, then $\mathsf{EXP} = \mathsf{NEXP}$. (Hint: padding argument.)
\end{exercise}

\begin{exercise}[Basic]
Show that $\mathsf{L} \subseteq \mathsf{NL} \subseteq \mathsf{P}$. Why don't we know if $\mathsf{L} = \mathsf{P}$?
\end{exercise}

\begin{exercise}[Basic]
Prove that $\mathsf{P} \subseteq \mathsf{NP} \cap \mathsf{coNP}$. Show that if $\mathsf{NP} \neq \mathsf{coNP}$, then $\mathsf{P} \neq \mathsf{NP}$.
\end{exercise}

\begin{exercise}[Basic]
Show that $\mathsf{BPP} \subseteq \mathsf{PSPACE}$ by simulating all random strings using depth-first search.
\end{exercise}

\begin{exercise}[Basic]
Prove that $\mathsf{coNP}$ is closed under polynomial-time reductions: if $A \le_p B$ and $B \in \mathsf{coNP}$, then $A \in \mathsf{coNP}$.
\end{exercise}

\begin{exercise}[Intermediate]
Prove Savitch's Theorem: $\mathsf{NSPACE}(s(n)) \subseteq \mathsf{DSPACE}(s(n)^2)$ for $s(n) \ge \log n$. (Hint: recursive reachability in configuration graph.)
\end{exercise}

\begin{exercise}[Intermediate]
Prove the padding argument: $\mathsf{P} = \mathsf{NP}$ implies $\mathsf{EXP} = \mathsf{NEXP}$. Show that the converse is false (padding cannot be reversed for smaller classes).
\end{exercise}

\begin{exercise}[Intermediate]
Prove Adleman's Theorem: $\mathsf{BPP} \subseteq \mathsf{P}/\mathsf{poly}$. (Hint: error reduction via majority vote, then union bound over all inputs.)
\end{exercise}

\begin{exercise}[Intermediate]
Show that $\mathsf{BPP} \subseteq \Sigma_2^{\mathsf{P}} \cap \Pi_2^{\mathsf{P}}$ (Sipser-G\'acs Theorem). This shows BPP lies low in the polynomial hierarchy.
\end{exercise}

\begin{exercise}[Intermediate]
Prove that the Parity function cannot be computed by a depth-2 $\mathsf{AC}^0$ circuit of subexponential size. Generalize this to depth $d$ using the switching lemma.
\end{exercise}

\begin{exercise}[Intermediate]
Prove that $k$-Vertex Cover is FPT via bounded search tree: branch on an edge, recurse with $k-1$, and analyze the branching factor.
\end{exercise}

\begin{exercise}[Intermediate]
Show that if $\mathsf{NP} = \mathsf{coNP}$, then the Polynomial Hierarchy collapses to $\mathsf{NP}$ ($\mathsf{PH} = \mathsf{NP}$).
\end{exercise}

\begin{exercise}[Challenge]
Research the \emph{Natural Proofs} barrier (Razborov-Rudich 1997). Explain why any proof that separates P from NP via circuit lower bounds must either be non-constructive or fail to apply to a ``natural'' property.
\end{exercise}

\begin{exercise}[Challenge]
Research \emph{fine-grained complexity}: What is the Orthogonal Vectors Hypothesis? How does it imply conditional lower bounds for Edit Distance, LCS, and Dynamic Programming problems?
\end{exercise}

\begin{exercise}[Challenge]
Research \emph{derandomization}: How does the Nisan-Wigderson generator work? Why does $\mathsf{E} \not\subseteq \mathsf{SIZE}(2^{\epsilon n})$ imply $\mathsf{P} = \mathsf{BPP}$?
\end{exercise}

\begin{exercise}[Challenge]
Prove the Resolution width lower bound: show that the pigeonhole principle $\mathsf{PHP}_n^{n+1}$ requires width $\Omega(n)$ in Resolution, and therefore size $\exp(\Omega(n))$.
\end{exercise}

\begin{exercise}[Challenge]
Research Impagliazzo's five worlds. For each world, determine: (1) what computational assumptions hold; (2) what cryptographic primitives exist; (3) what is the status of P vs NP.
\end{exercise}

\begin{exercise}[Challenge]
Show that $k$-Clique is $\mathsf{W}[1]$-complete via a parameterized reduction from $k$-Independent Set. Explain why this rules out FPT algorithms for Clique (unless $\mathsf{FPT} = \mathsf{W}[1]$).
\end{exercise}

\begin{exercise}[Challenge]
Research the PCP Theorem and explain its statement. Why does it imply that MAX-3SAT is NP-hard to approximate within a factor of $7/8 + \epsilon$?
\end{exercise}

\begin{exercise}[Computational]
Implement a CDCL SAT solver (or use PySAT/Z3). Test on random 3-SAT at the phase transition ($m/n \approx 4.26$). Plot solving time vs. $n$. Compare with DPLL.
\end{exercise}

\begin{exercise}[Computational]
Implement a simple parameterized algorithm for Vertex Cover (branching on edges). Measure runtime vs. $k$ and $n$. Compare with ILP formulation.
\end{exercise}

\begin{exercise}[Computational]
Implement a DP algorithm for Edit Distance. Then try to find sub-quadratic speedups for special cases (e.g., when edit distance is small). Compare with the fine-grained lower bound.
\end{exercise}

%======================================================================
% TIMELINE
%======================================================================
\section{Timeline}
\label{sec:complexity:timeline}
\addcontentsline{toc}{section}{Timeline}

\begin{tabularx}{\linewidth}{|l|X|}
\hline
\textbf{Year} & \textbf{Event} \\ \hline
1956 & G\"odel's letter to von Neumann (first P vs NP intuition) \\ \hline
1960 & Rabin: Degree of difficulty \\ \hline
1965 & Hartmanis-Stearns: Complexity classes, hierarchy theorems \\ \hline
1965 & Edmonds: Polynomial time = tractable (``good algorithms'') \\
 \hline
1965 & Cobham: Intrinsic computational difficulty \\
 \hline
1967 & Paterson: Time-bounded TMs \\
 \hline
1970 & Savitch: NSPACE($s$) $\subseteq$ DSPACE($s^2$) \\
 \hline
1971 & Cook: SAT is NP-complete \\
 \hline
1972 & Karp: 21 NP-complete problems \\
 \hline
1973 & Levin: Universal search problems (USSR) \\
 \hline
1974 & Fagin: Descriptive complexity (NP = $\exists$SO) \\
 \hline
1975 & Ladner: NP-intermediate problems \\
 \hline
1975 & Baker-Gill-Solovay: Relativization (P vs NP oracles) \\
 \hline
1976 & Pippenger-Fischer: Oblivious TM simulation \\
 \hline
1978 & Adleman: BPP $\subseteq$ P/poly \\
 \hline
1979 & Cook-Reckhow: Proof complexity \\
 \hline
1983 & Yao: BPP $\subseteq$ P/poly (alternative proof); Yao's XOR lemma \\
 \hline
1984 & Furst-Saxe-Sipser, Ajtai: Parity $\notin$ AC$^0$ \\
 \hline
1986 & H{\aa}stad: Tight switching lemma, optimal AC$^0$ lower bounds \\
 \hline
1987 & Razborov-Smolensky: AC$^0$[p] lower bounds \\
 \hline
1991 & Babai-Moran: Arthur-Merlin games \\
 \hline
1994 & Shor: Factoring in BQP \\
 \hline
1994 & Nisan-Wigderson: Pseudorandom generator from hardness \\
 \hline
1996 & Grover: Quantum search ($O(\sqrt{N})$) \\
 \hline
1997 & Razborov-Rudich: Natural Proofs barrier \\
 \hline
1997 & Impagliazzo-Wigderson: P = BPP from circuit lower bounds \\
 \hline
1999 & Downey-Fellows: Parameterized complexity (W-hierarchy) \\
 \hline
2001 & Dinur: PCP theorem (alternative proof) \\
 \hline
2001 & Ben-Sasson-Wigderson: Size-width tradeoff for Resolution \\
 \hline
2008 & Aaronson-Wigderson: Algebrization barrier \\
 \hline
2010s & Fine-grained complexity, optimal SETH-based lower bounds \\
 \hline
2014 & Williams: ACC lower bounds via circuit satisfiability \\
 \hline
2019 & Google Sycamore: Quantum supremacy experiment \\
 \hline
2020s & Meta-complexity: MCSP, worst-case to average-case reductions \\
 \hline
\end{tabularx}

%======================================================================
% CITATIONS
%======================================================================
\section{Citations}
\label{sec:complexity:citations}

\begin{enumerate}
  \item \cite{HartmanisStearns1965}
  \item \cite{Edmonds1965}
  \item \cite{Cobham1965}
  \item \cite{Cook1971}
  \item \cite{Karp1972}
  \item \cite{Levin1973}
  \item \cite{Savitch1970}
  \item \cite{Ladner1975}
  \item \cite{BakerGillSolovay1975}
  \item \cite{FurstSaxeSipser1984}
  \item \cite{Ajtai1983}
  \item \cite{RazborovSmolensky1987}
  \item \cite{RazborovRudich1997}
  \item \cite{NisanWigderson1994}
  \item \cite{ImpagliazzoWigderson1997}
  \item \cite{Shor1994}
  \item \cite{Grover1996}
  \item \cite{DowneyFellows1999}
  \item \cite{CookReckhow1979}
  \item \cite{BenSassonWigderson2001}
  \item \cite{arorabarack2009}
  \item \cite{goldreich2004}
  \item \cite{Stockmeyer1976}
  \item \cite{Valiant1984}
  \item \cite{Fagin1974}
  \item \cite{Williams2014}
\end{enumerate}

%======================================================================
% INDEX ENTRIES
%======================================================================
\index{computational complexity}
\index{P vs NP}
\index{polynomial time}
\index{NP-complete}
\index{hierarchy theorems}
\index{Hartmanis, Juris}
\index{Stearns, Richard}
\index{Edmonds, Jack}
\index{Cobham, Alan}
\index{Cook, Stephen}
\index{Karp, Richard}
\index{Levin, Leonid}
\index{Savitch's theorem}
\index{Extended Church-Turing Thesis}
\index{fine-grained complexity}
\index{parameterized complexity}
\index{natural proofs}
\index{circuit complexity}
\index{derandomization}
\index{pseudorandomness}
\index{BQP}
\index{quantum complexity}
\index{meta-complexity}
\index{proof complexity}
\index{Relativization}
\index{Switching Lemma}
\index{Polynomial Hierarchy}
\index{FPT}
\index{W-hierarchy}
\index{Resolution}
\index{Cutting Planes}
\index{Sum-of-Squares}
\index{SETH}
\index{OV Hypothesis}
\index{3SUM}
\index{APSP}
\index{kernelization}
\index{treewidth}
\index{Impagliazzo's Five Worlds}
\index{PAC learning}
\index{descriptive complexity}
